\section{Tag 1 — Prüfziffern und Fehlererkennung}

\subsection*{Bleistiftübung 1 — Wiederholungscode}

\begin{enumerate}
\item Mehrheitsentscheid: \texttt{111}$\to 1$, \texttt{010}$\to 0$,
  \texttt{001}$\to 0$, \texttt{110}$\to 1$, \texttt{000}$\to 0$,
  \texttt{101}$\to 1$.
\item Pro Block korrigieren: 1 Bitfehler. Erkennen (ohne zu
  korrigieren): 2 Bitfehler. (3 Fehler = alle gekippt → wird falsch
  korrigiert!)
\item Redundanz: 2 von 3 Bits sind Redundanz, also etwa 67\,\%.
\item Viel zu teuer. In der Praxis braucht man Verfahren, die mit
  z.\,B.\ 10\,\% Redundanz auskommen.
\end{enumerate}

\subsection*{Bleistiftübung 2 — Modulo}

\begin{center}
\begin{tabular}{lc}
  \toprule
  Aufgabe & Ergebnis \\
  \midrule
  $29 \bmod 7$         & 1 \\
  $100 \bmod 9$        & 1 \\
  $1000 \bmod 11$      & 10 \\
  $(47 + 38) \bmod 10$ & 5 \\
  $(8 \cdot 7) \bmod 10$ & 6 \\
  $0 \bmod 13$         & 0 \\
  \bottomrule
\end{tabular}
\end{center}

\subsection*{Bleistiftübung 3 — Parität}

(a)
\begin{center}
\begin{tabular}{ccc}
  \toprule
  Datenwort & P & Gesendet \\
  \midrule
  1010101 & 0 & 10101010 \\
  1111000 & 0 & 11110000 \\
  0000001 & 1 & 00000011 \\
  1100110 & 0 & 11001100 \\
  0000000 & 0 & 00000000 \\
  \bottomrule
\end{tabular}
\end{center}

(b) Anzahl der Einsen zählen, ungerade = fehlerhaft:
\begin{itemize}
\item \texttt{10110011} → 5 Einsen → \textbf{fehlerhaft}
\item \texttt{11000001} → 3 Einsen → \textbf{fehlerhaft}
\item \texttt{00000000} → 0 → ok
\item \texttt{01010100} → 3 → \textbf{fehlerhaft}
\item \texttt{11111111} → 8 → ok
\end{itemize}

(c) Bei zwei gekippten Bits kippt auch die Anzahl der Einsen zweimal um
→ die Parität bleibt \textbf{gleich} → Fehler \textbf{nicht erkannt}.

(d) Parität erkennt genau \textbf{eine ungerade Anzahl gekippter Bits}
(1, 3, 5, \dots Fehler). Gerade Fehleranzahlen bleiben unbemerkt.

\subsection*{Bleistiftübung 4 — EAN-13}

(b) \texttt{400638133393}:
\begin{verbatim}
4·1 + 0·3 + 0·1 + 6·3 + 3·1 + 8·3 + 1·1 + 3·3 + 3·1 + 3·3 + 9·1 + 3·3
= 4 + 0 + 0 + 18 + 3 + 24 + 1 + 9 + 3 + 9 + 9 + 9 = 89
(10 − 89 mod 10) mod 10 = (10 − 9) mod 10 = 1
\end{verbatim}
→ Prüfziffer \textbf{1}, voller Code: \texttt{4006381333931}.

(c) Mit dem Verfahren ausrechnen:
\begin{itemize}
\item \texttt{4006381333931} → gültig (Prüfziffer 1 \checkmark)
\item \texttt{9783446435001} → gültig (ISBN-13, Prüfziffer 1
  \checkmark)
\item \texttt{4012345678902} → Summe
  $= 4 \cdot 1 + 0 \cdot 3 + 1 \cdot 1 + 2 \cdot 3 + 3 \cdot 1 + 4 \cdot 3 + 5 \cdot 1 + 6 \cdot 3 + 7 \cdot 1 + 8 \cdot 3 + 9 \cdot 1 + 0 \cdot 3
  = 4+0+1+6+3+12+5+18+7+24+9+0 = \mathbf{89}$, Prüfziffer also
  \textbf{1}, nicht 2 → \textbf{gefälscht}.
\end{itemize}

\subsection*{Aufgabe 3.1 — Lösung}

\pythoncodeteil{tag1/paritaet.py}{test-einzel}

Ergebnis: 1024/1024 (100\,\%).

\subsection*{Aufgabe 3.2 — Lösung}

\pythoncodeteil{tag1/paritaet.py}{test-doppel}

Ergebnis: 0/3584 (0\,\%!) — Parität erkennt KEINEN Doppelfehler.

\subsection*{Aufgaben 5.1 / 5.2 — Lösung}

Bei \texttt{4006381333931} werden Vertauschungen an Position 1↔2
(\texttt{00}/\texttt{06} — egal, weil eine Ziffer 0 ist) usw. Wenn
man systematisch alle EANs ausprobiert, erkennt man:

\begin{quote}
\textbf{Ein Zifferndreher zwischen benachbarten Ziffern $a$ und $b$
wird genau dann nicht erkannt, wenn $|a - b| = 5$.}
\end{quote}

Begründung: die Differenz, die der Tausch in der Summe macht, ist
$(b \cdot 1 + a \cdot 3) - (a \cdot 1 + b \cdot 3) = 2(a - b)$. Damit
das modulo 10 verschwindet, muss $2(a-b) \equiv 0 \pmod{10}$, also
$a - b \equiv 0 \pmod{5}$. Bei einer echten Vertauschung ($a \ne b$)
heißt das: $|a-b| = 5$.

→ Von den 45 möglichen Ziffernpaaren mit $a \ne b$ haben 5 Paare die
Differenz 5: $(0,5)$, $(1,6)$, $(2,7)$, $(3,8)$, $(4,9)$. Das sind
$5/45 \approx 11\,\%$ aller möglichen Zifferndreher, die EAN-13
\textbf{nicht} erkennt.

\subsection*{Aufgabe 5.3 — Lösung}

\pythoncodeteil{tag1/ean13.py}{einzelfehler}

Alle 117 Einzelfehler werden erkannt.

\subsection*{Antworten zu den drei Reflexionsfragen}

\begin{enumerate}
\item \textbf{Erkennen} heißt nur „etwas stimmt nicht“. \textbf{Korrigieren}
  heißt „ich weiß, was ursprünglich da stand“. Dafür müssen sich
  gültige Codewörter so stark unterscheiden, dass sich nach einem
  Fehler immer noch das ursprüngliche Codewort eindeutig „am nächsten“
  anbietet → führt zur \textbf{Hamming-Distanz} (Tag 2/3).
\item Statt der Gewichte 1 und 3 könnte man andere Gewichte (z.\,B.\
  mit Modulo 11) verwenden — das macht die ISBN-10 (kommt morgen).
\item Spoiler: für 1-Bit-Korrektur reichen erstaunlich wenige
  Zusatzbits — siehe \textbf{Hamming-Code (7,4)} in Tag 3.
\end{enumerate}
